% Part: first-order-logic
% Chapter: syntax-and-semantics
% Section: first-order-languages

\documentclass[../../../include/open-logic-section]{subfiles}

\begin{document}

\olfileid{fol}{syn}{fol}

\olsection{પ્રથમ-ક્રમ ભાષાઓ}


પ્રથમ-ક્રમ તર્કશાસ્ત્રની અભિવ્યક્તિઓ, \emph{!!{variable}s},
\emph{!!{constant}s}, \emph{!!{predicate}s} અને ક્યારેક
\emph{!!{function}s} ધરાવતા પાયાના સંકેતભંડોળમાંથી રચાય છે. આ
સંકેતોમાંથી તાર્કિક સંયોજકો, પરિમાણકો અને કૌંસ તથા અલ્પવિરામ જેવાં
વિરામચિહ્નો સાથે મળીને \emph{પદો} અને \emph{!!{formula}s} રચાય છે.

\begin{explain}
અનૌપચારિક રીતે, !!{predicate}s ગુણધર્મો અને સંબંધોનાં નામ છે,
!!{constant}s વ્યક્તિગત વસ્તુઓનાં નામ છે અને !!{function}s
પ્રતિચિત્રણોનાં નામ છે. !!{identity}~$\eq$ને બાદ કરતાં આ બધાં
\emph{અતાર્કિક સંકેતો} છે અને તે સાથે મળીને ભાષા બનાવે છે. કોઈ પણ
પ્રથમ-ક્રમ ભાષા~$\Lang L$ તેના અતાર્કિક સંકેતોથી નિર્ધારિત થાય છે.
સૌથી સામાન્ય કિસ્સામાં, $\Lang L$માં દરેક પ્રકારના અનંત સંખ્યામાં
સંકેતો હોય છે.
\end{explain}

સામાન્ય કિસ્સામાં પ્રથમ-ક્રમ તર્કશાસ્ત્રમાં આપણે નીચેના સંકેતો
વાપરીએ છીએ:

\begin{enumerate}
\item તાર્કિક સંકેતો
\begin{enumerate}
\item તાર્કિક સંયોજકો:
  \startycommalist
  \iftag{prvNot}{\ycomma $\lnot$ (નિષેધ)}{}%
  \iftag{prvAnd}{\ycomma $\land$ (સંયોજન)}{}%
  \iftag{prvOr}{\ycomma $\lor$ (વિયોજન)}{}%
  \iftag{prvIf}{\ycomma $\lif$ (!!{conditional})}{}%
  \iftag{prvIff}{\ycomma $\liff$ (!!{biconditional})}{}%
  \iftag{prvAll}{\ycomma $\lforall$ (સાર્વત્રિક પરિમાણક)}{}%
  \iftag{prvEx}{\ycomma $\lexists$ (અસ્તિત્વલક્ષી પરિમાણક)}{}.
\tagitem{prvFalse}{!!{falsity} માટેનો વિધાનાત્મક અચળ~$\lfalse$.}{}
\tagitem{prvTrue}{!!{truth} માટેનો વિધાનાત્મક અચળ~$\ltrue$.}{}
\item દ્વિસ્થાની !!{identity}~$\eq$.
\item !!{variable}sનો !!{denumerable}s ગણ: $\Obj v_0$, $\Obj v_1$,
  $\Obj v_2$, \dots
\end{enumerate}
\item પ્રથમ-ક્રમ તર્કશાસ્ત્રની \emph{પ્રમાણભૂત ભાષા} બનાવતા
  અતાર્કિક સંકેતો
\begin{enumerate}
\item દરેક $n>0$ માટે $n$-સ્થાની !!{predicate}sનો !!{denumerable}s
  ગણ: $\Obj A^n_0$, $\Obj A^n_1$, $\Obj A^n_2$, \dots
\item !!{constant}sનો !!{denumerable}s ગણ: $\Obj c_0$, $\Obj c_1$,
  $\Obj c_2$, \dots.
\item દરેક $n>0$ માટે $n$-સ્થાની !!{function}sનો !!{denumerable}s
  ગણ: $\Obj f^n_0$, $\Obj f^n_1$, $\Obj f^n_2$, \dots
\end{enumerate}
\item વિરામચિહ્નો: (, ) અને અલ્પવિરામ.
\end{enumerate}

આપણી મોટા ભાગની વ્યાખ્યાઓ અને પરિણામો પ્રથમ-ક્રમ તર્કશાસ્ત્રની
સંપૂર્ણ પ્રમાણભૂત ભાષા માટે રજૂ થશે. જોકે ઉપયોગ અનુસાર આપણે ભાષાને
માત્ર થોડાંક !!{predicate}s, !!{constant}s અને !!{function}s સુધી
મર્યાદિત પણ કરી શકીએ છીએ.

\begin{ex}
અંકગણિતની ભાષા~$\Lang L_A$માં એક દ્વિસ્થાની !!{predicate}~$<$, એક
!!{constant}~$\Obj 0$, એક એકસ્થાની !!{function}~$\prime$ અને બે
દ્વિસ્થાની !!{function}s~$+$ તથા~$\times$ હોય છે.
\end{ex}

\begin{ex}
ગણસિદ્ધાંતની ભાષા~$\Lang L_Z$માં માત્ર એક દ્વિસ્થાની
!!{predicate}~$\in$ હોય છે.
\end{ex}

\begin{ex}
ક્રમોની ભાષા~$\Lang L_\le$માં માત્ર દ્વિસ્થાની
!!{predicate}~$\le$ હોય છે.
\end{ex}

ફરી યાદ રાખો કે આ પરંપરાઓ છે: સત્તાવાર રીતે આ માત્ર ઉપનામો છે;
દાખલા તરીકે, $<$, $\in$ અને $\le$ એ $\Obj A^2_0$નાં ઉપનામો છે,
$\Obj 0$ એ $\Obj c_0$નું, $\prime$ એ $\Obj f^1_0$નું, $+$ એ
$\Obj f^2_0$નું અને $\times$ એ $\Obj f^2_1$નું ઉપનામ છે.

\iftag{defNot,defOr,defAnd,defIf,defIff,defTrue,defFalse,defEx,defAll}{%
ઉપર રજૂ કરેલા મૂળભૂત સંયોજકો અને
\iftag{notprvEx,notprvAll}{પરિમાણક}{પરિમાણકો} ઉપરાંત આપણે નીચેના
\emph{વ્યાખ્યાયિત} સંકેતો પણ વાપરીએ છીએ:
\startycommalist
  \iftag{defNot}{\ycomma $\lnot$ (નિષેધ)}{}%
  \iftag{defAnd}{\ycomma $\land$ (સંયોજન)}{}%
  \iftag{defOr}{\ycomma $\lor$ (વિયોજન)}{}%
  \iftag{defIf}{\ycomma $\lif$ (!!{conditional})}{}%
  \iftag{defIff}{\ycomma $\liff$ (!!{biconditional})}{}%
  \iftag{defAll}{\ycomma $\lforall$ (સાર્વત્રિક પરિમાણક)}{}%
  \iftag{defEx}{\ycomma $\lexists$ (અસ્તિત્વલક્ષી પરિમાણક)}{}%
  \iftag{defFalse}{\ycomma !!{falsity}~$\lfalse$}{}%
  \iftag{defTrue}{\ycomma !!{truth}~$\ltrue$}{}.
}{}
\sourcecorrection{OLSYN-001}{મૂળની
\texttt{first-order-languages.tex}માં વ્યાખ્યાયિત $\ltrue$ માટેની
અંદરની \texttt{\string\iftag} આજ્ઞાનો ત્રીજો દલીલખંડ છૂટી ગયો છે અને
તેના સ્થાને આવેલું અંતકૌંસ બહારની \texttt{\string\iftag} આજ્ઞાનો
બીજો દલીલખંડ અકાળે બંધ કરે છે. અહીં બંને આજ્ઞાઓનું સંતુલિત,
ત્રણ-દલીલવાળું રૂપ પુનઃસ્થાપિત કર્યું છે.}

\begin{tagblock}{defNot,defOr,defAnd,defIf,defIff,defTrue,defFalse,defEx,defAll}
\begin{explain}
વ્યાખ્યાયિત સંકેત સત્તાવાર રીતે ભાષાનો ભાગ નથી, પણ અનૌપચારિક
સંક્ષેપ તરીકે રજૂ થાય છે: માત્ર મૂળભૂત સંકેતો વાપરીએ તો ઘણાં લાંબાં
થતાં સૂત્રોને તેના વડે ટૂંકાં લખી શકાય છે. આ દેખીતો લાભ છે. જોકે વધુ
મોટો લાભ એ છે કે સાબિતીઓ ટૂંકી થાય છે. કોઈ સંકેત મૂળભૂત હોય તો
સાબિતીઓમાં તેને અલગથી ધ્યાનમાં લેવો પડે. તેથી મૂળભૂત સંકેતો જેટલા
વધુ, આપણી સાબિતીઓ એટલી લાંબી.
\end{explain}
\end{tagblock}

% Alternate symbols

\begin{intro}
ઉપર આપણે વાપર્યાં તેનાથી જુદી પરિભાષા અને સંકેતોથી તમે પરિચિત હોઈ
શકો. તર્કશાસ્ત્રનાં પુસ્તકો (અને શિક્ષકો) ``નિષેધ'' માટે સામાન્ય રીતે
$\sim$, $\neg$ અથવા~! અને ``સંયોજન'' માટે $\wedge$, $\cdot$ અથવા
$\&$ વાપરે છે. ``શરતી વિધાન'' અથવા ``પ્રેરણ'' માટે સામાન્ય રીતે
$\rightarrow$, $\Rightarrow$ અને $\supset$ વપરાય છે.
\iftag{prvIff,defIff}{``દ્વિશરતી વિધાન'', ``દ્વિમુખી પ્રેરણ'' અથવા
  ``(સત્યમૂલ્યલક્ષી) સમતુલ્યતા'' માટેના સંકેતો $\leftrightarrow$,
  $\Leftrightarrow$ અને $\equiv$ છે.}{}
\iftag{prvFalse,defFalse}{$\lfalse$ સંકેતને જુદી જગ્યાએ ``અસત્યતા'',
  ``ફાલ્સમ'', ``અસંગતિ'' અથવા ``તળિયું'' કહે છે.}{}
\iftag{prvTrue,defTrue}{$\ltrue$ સંકેતને જુદી જગ્યાએ ``સત્ય'',
  ``વેરમ'' અથવા ``ટોચ'' કહે છે.}{}

લેટિન મૂળાક્ષરના આરંભનાં નાનાં અક્ષરો (દા.ત., $a$, $b$, $c$)
!!{constant}s માટે (ક્યારેક તેમને નામો કહેવાય છે) અને અંતનાં નાનાં
અક્ષરો (દા.ત., $x$, $y$, $z$) !!{variable}s માટે વાપરવાની પરંપરા
છે. પરિમાણકો !!{variable}s સાથે જોડાય છે, જેમ કે~$x$; સાર્વત્રિક
પરિમાણક માટે $\forall x$, $(\forall x)$, $(x)$, $\Pi x$,
$\bigwedge_x$ અને અસ્તિત્વલક્ષી પરિમાણક માટે $\exists x$,
$(\exists x)$, $(Ex)$, $\Sigma x$, $\bigvee_x$ જેવી સંકેત-ભિન્નતાઓ
વપરાય છે.
\end{intro}

\begin{explain}
આપણે બધા વિધાનાત્મક કારકો અને બંને પરિમાણકોને ભાષાના મૂળભૂત સંકેતો
ગણી શકીએ. તેના બદલે મૂળભૂત સંકેતોનો નાનો ભંડાર પસંદ કરીને બીજા
!!{operator}sને વ્યાખ્યાયિત પણ ગણી શકીએ. બૂલિયન કારકોના
``સત્યતાફલનલક્ષી રીતે પૂર્ણ'' ગણોમાં $\{ \lnot, \lor \}$,
$\{ \lnot, \land \}$ અને $\{ \lnot, \lif\}$નો સમાવેશ થાય છે;
આમાંના કોઈ પણ ગણને કોઈ એક પરિમાણક સાથે જોડીને અભિવ્યક્તિક્ષમ રીતે
પૂર્ણ પ્રથમ-ક્રમ ભાષા મળે છે.

તમે બીજા બે !!{operator}sથી પણ પરિચિત હોઈ શકો: હેન્રી શેફરના નામ
પરથી ઓળખાતો શેફર સ્ટ્રોક~$|$ અને ક્વાઇનના ડૅગર તરીકે પણ ઓળખાતો
પિયર્સનો તીર~$\downarrow$. અનુક્રમે ``નૅન્ડ'' અને ``નૉર''ના તેમના
સામાન્ય વાંચન હેઠળ આ દરેક કારક પોતે જ સત્યતાફલનલક્ષી રીતે પૂર્ણ છે.
\end{explain}

\end{document}
